% !TEX root = ../demo.tex
% =================================================================
%  templates/lists-and-math.tex
%  Текстовые и математические конструкции конспекта по матанализу
% =================================================================

\subsection{Списки}

\begin{itemize}
	\item Ограниченность: $\|\mathcal{A}\| < +\infty$.
	\item Непрерывность в точке.
	\item Равномерная непрерывность.
\end{itemize}

\begin{enumerate}
	\item Зафиксируем $\eps > 0$.
	\item Подберём $\delta > 0$ по определению непрерывности.
	\item Оценим приращение и перейдём к пределу.
\end{enumerate}

\subsection{Выделения в тексте}

\textbf{Жирным} вводится термин; \emph{курсивом} --- уточняющий акцент;
\texttt{моноширинным} --- имена команд. Тире набирается \verb|---| (---),
кавычки --- \verb|<<...>>| (<<так>>), а тире после термина в эталоне
идёт как \verb|~---| (без разрыва): пусть $X$~--- банахово пространство.

\subsection{Кванторы, пределы, оценки}

\[
\lim_{\substack{y\to x\\ y\in\Omega}} f(y) = f(x),
\qquad
\forall \eps>0\ \exists \delta>0 \st |y-x|<\delta \Rightarrow |f(y)-f(x)|<\eps.
\]

Двусторонняя оценка --- типовая конструкция матанализа:
\[
|a_k| \le \sqrt{\sum_{j=1}^{n} a_j^2} = |a| \le \sqrt{n}\,\max_{1\le j\le n}|a_j|.
\]

Цепочка неравенств с интегралами:
\[
\int_{n-1}^{n}\ln x\,dx \;<\; \ln n \;<\; \int_{n}^{n+1}\ln x\,dx .
\]

\subsection{Нормы и скалярные произведения}

Макросы из \texttt{style/math-commands.tex}:
\[
\norma{x},\qquad \abs{x},\qquad \scal{x}{y},\qquad
\Lp{f},\qquad \Lq{g},\qquad \norm{f}{2}.
\]

\subsection{Сокращение дробей: пакет \texttt{cancel}}

\[
\frac{\cancel{(x-a)}\,g(x)}{\cancel{(x-a)}} = g(x),
\qquad
\cancelto{0}{\frac{1}{n}} \ \text{при } n\to\infty.
\]

\subsection{Разбор случаев и системы}

\[
f(x) =
\begin{cases}
	x\sin\frac1x, & x \neq 0,\\
	0,            & x = 0;
\end{cases}
\qquad
\begin{cases}
	\|\mathcal{A}x\|_Y \le C\|x\|_X,\\
	C \ge 0 .
\end{cases}
\]

\subsection{Многострочная выкладка}

\begin{align*}
	\Bigl\| x - \sum_{k=1}^{N}\lambda_k e_k \Bigr\|^2
	&= \|x\|^2 - 2\sum_{k=1}^{N}\lambda_k c_k + \sum_{k=1}^{N}\lambda_k^2 \\
	&= \|x\|^2 - \sum_{k=1}^{N} c_k^2 + \sum_{k=1}^{N}(\lambda_k - c_k)^2 .
\end{align*}

\subsection{Ссылки}

На бокс: \verb|\ref{thm:demo-thm}| $\to$ \ref{thm:demo-thm}.
На другой билет --- просто текстом: см. вопрос~38.

\subsection{Хайлайты}

Шесть цветов палитры из \texttt{style/highlight.tex} и короткие команды к ним:
\hlY{жёлтый}, \hlB{синий}, \hlG{зелёный}, \hlR{красный},
\hlO{оранжевый}, \hlV{фиолетовый}. Произвольный цвет палитры ---
\verb|\hlc{hlBlue}{текст}|.

Те же команды работают внутри формул:
\[
c_k = \hlR{\sum_{t=0}^{n-1} a_t\, b_{k-t}}, \qquad
\hlY{k = 0,\dots,2n-2}.
\]

Заливка рассчитана на короткий фрагмент и \emph{не переносится} на следующую
строку: длинный кусок либо разбивают на несколько \verb|\hl...|, либо
выделяют боксом.

\subsection{Помеченная формула и ссылка на неё}

Нумерованные формулы нумеруются внутри \verb|\section| --- как боксы и
рисунки. Метка и ссылка:

\begin{equation}\label{eq:demo}
	\pder{u}{x} + \pder{u}{y}\, f(x,y) \equiv 0 .
\end{equation}

Ссылка в тексте: критерий~\eqref{eq:demo}. Для формул без номера
по-прежнему используется \verb|\[ ... \]|.
